\documentclass[light,minted=false]{cuzbeamer} % tema claro, sin resaltado de código

\usepackage{graphicx}

\title{Aplicación de Inteligencia Artificial Generativa en la Visualización de Datos}
\subtitle{Con Python, Plotly y Dash}
\author{Dr. Alejandro González Turrubiates}
\institute{Universidad Autónoma de Tamaulipas}
\date{2 al 4 de septiembre de 2026}
\email{agturrubiates@uat.edu.mx}

% Referencia textual del congreso donde se imparte el taller (el logo del
% evento ya ocupa, en su lugar, la esquina superior derecha de cada slide:
% ver images/logo_SIP_UAT-trimmed.png y images/logo_SIP_UAT-negro.png).
\AtBeginDocument{%
    \addtobeamertemplate{title page}{}{%
        \begin{tikzpicture}[remember picture,overlay]
            \node[anchor=south west,xshift=0.3cm,yshift=0.35cm,align=left,font=\tiny,text=white]
                at (current page.south west)
                {En el marco del Congreso Internacional CIREDII 2026 --- Colima};
        \end{tikzpicture}%
    }%
}

\begin{document}
    \maketitle

    \begin{frame}{Objetivos del taller}
        \begin{itemize}
            \item Formular prompts efectivos para generar visualizaciones
            \item Generar y ajustar gráficas a partir de datos reales
            \item \textbf{Verificar} los resultados generados por IA
            \item Aplicar buenas prácticas de uso responsable de la IA
        \end{itemize}
    \end{frame}

    \begin{frame}{Agenda del día (9:00--14:00)}
        \begin{enumerate}
            \item Fundamentos de prompting
            \item Plotly guiado
            \item Práctica libre
            \item De Plotly a Dash
            \item Buenas prácticas y cierre
        \end{enumerate}
    \end{frame}

    \begin{frame}{Antes de empezar}
        Verifica que tienes instalado:
        \begin{itemize}
            \item Python 3.11+
            \item Una CLI agentic (Gemini CLI o Claude Code)
        \end{itemize}

        \vspace{1em}
        Ver \texttt{requisitos/} si algo falta.
    \end{frame}

    % =========================================================
    % BLOQUE 1 — Fundamentos de prompting (9:20--9:50)
    % =========================================================
    \section{Fundamentos de prompting}

    \begin{frame}{¿Qué es un agente de IA en terminal?}
        Un programa que corre en tu terminal y puede:
        \begin{itemize}
            \item \textbf{Leer} tu instrucción en lenguaje natural
            \item \textbf{Escribir} código Python por ti
            \item \textbf{Ejecutarlo} y mostrarte el resultado
            \item \textbf{Corregirlo} si algo sale mal, si se lo pides
        \end{itemize}

        \vspace{1em}
        Tú nunca escribes el código directamente — lo pides, lo revisas y lo verificas.
    \end{frame}

    \begin{frame}{El patrón de trabajo}
        \begin{enumerate}
            \item \textbf{Prompt} — describes qué gráfica necesitas y con qué datos
            \item \textbf{Código} — la IA escribe un script Python
            \item \textbf{Gráfica} — el script corre y produce un resultado visual
            \item \textbf{Verificación} — revisas que el resultado sea correcto
        \end{enumerate}

        \vspace{1em}
        Este ciclo se repite las veces que haga falta, ajustando el prompt.
    \end{frame}

    \begin{frame}{Cómo pedir un buen prompt}
        \begin{itemize}
            \item Sé específico: qué tipo de gráfica, con qué columnas, de qué archivo
            \item Pide un paso a la vez, no una app completa de un jalón
            \item Pide que te explique en una oración qué hizo, antes de correrlo
            \item Si algo está mal, describe \textbf{qué} está mal — no pidas "hazlo de nuevo"
        \end{itemize}
    \end{frame}

    % =========================================================
    % Checkpoint de instalación (9:50--10:00)
    % =========================================================
    \begin{frame}{Checkpoint de instalación}
        Antes de la práctica, confirma que tienes:
        \begin{itemize}
            \item Python 3.11+ funcionando (\texttt{python --version})
            \item Tu CLI agentic autenticada y respondiendo
        \end{itemize}

        \vspace{1em}
        Guías completas en \texttt{requisitos/} si necesitas resolver algo ahora.
    \end{frame}

    % =========================================================
    % BLOQUE 2 — Plotly guiado (10:00--11:00)
    % =========================================================
    \section{Plotly guiado}

    \begin{frame}{Objetivo del bloque}
        Generar tus primeras gráficas con Plotly, guiado paso a paso, usando
        datos abiertos reales.

        \vspace{1em}
        \textbf{Tú} rediges los prompts. \textbf{La IA} escribe y corre el código.
        \textbf{Ambos} verifican el resultado.
    \end{frame}

    \begin{frame}{Qué vamos a construir}
        \begin{itemize}
            \item Gráfica de \textbf{barras} — comparar categorías
            \item Gráfica de \textbf{líneas} — ver una tendencia en el tiempo
            \item \textbf{Mapa} — ver un dato distribuido geográficamente
        \end{itemize}
    \end{frame}

    \begin{frame}{De dónde salen los datos}
        Vamos a trabajar con \textbf{datos abiertos} (INEGI y fuentes públicas de
        México) — datos reales, no de ejemplo.

        \vspace{1em}
        Ver \texttt{datos/} para las fuentes y los archivos específicos de este bloque.
    \end{frame}

    % =========================================================
    % Descanso (11:00--11:15)
    % =========================================================
    \begin{standout}[Descanso]
        11:00 -- 11:15
    \end{standout}

    % =========================================================
    % BLOQUE 3 — Práctica libre (11:15--12:15)
    % =========================================================
    \section{Práctica libre}

    \begin{frame}{Instrucciones}
        \begin{enumerate}
            \item Elige una pregunta sobre los datos disponibles en \texttt{datos/}
            \item Pídele a tu CLI agentic una gráfica que la responda
            \item Ajusta el prompt hasta que el resultado te convenza
            \item Verifica antes de darlo por bueno
        \end{enumerate}
    \end{frame}

    \begin{frame}{Checklist de verificación}
        Antes de dar por buena una gráfica generada por IA, revisa:
        \begin{itemize}
            \item ¿Los ejes muestran lo que dicen mostrar (unidades, escala)?
            \item ¿Los números coinciden con lo que esperarías del dataset?
            \item ¿La IA usó el archivo real, o inventó valores de ejemplo?
        \end{itemize}
    \end{frame}

    % =========================================================
    % Descanso corto (12:15--12:30)
    % =========================================================
    \begin{standout}[Descanso corto]
        12:15 -- 12:30
    \end{standout}

    % =========================================================
    % BLOQUE 4 — De Plotly a Dash (12:30--13:30)
    % =========================================================
    \section{De Plotly a Dash}

    \begin{frame}{¿Qué es Dash?}
        Un framework de Python que sirve tus gráficas de Plotly como una
        \textbf{aplicación web interactiva}, corriendo en un puerto local
        (\texttt{localhost}).

        \vspace{1em}
        Plotly genera la gráfica; Dash la convierte en dashboard.
    \end{frame}

    \begin{frame}{Qué vamos a construir}
        Un dashboard que reúne las gráficas de los bloques anteriores en una
        sola página, con al menos un control interactivo (por ejemplo, un
        filtro) que actualiza la gráfica.
    \end{frame}

    \begin{frame}{Cómo lo vamos a hacer}
        \begin{enumerate}
            \item Partimos del código Plotly que ya generamos
            \item Le pedimos a la IA que lo integre en una app Dash
            \item Corremos la app y la abrimos en el navegador (\texttt{localhost})
            \item Verificamos que el control interactivo funcione
        \end{enumerate}
    \end{frame}

    % =========================================================
    % BLOQUE 5 — Buenas prácticas y cierre (13:30--14:00)
    % =========================================================
    \section{Buenas prácticas y cierre}

    \begin{frame}{Buenas prácticas de uso responsable}
        \begin{itemize}
            \item Verifica siempre — la IA puede inventar datos con total confianza
            \item No compartas datos sensibles o confidenciales con herramientas de IA
            \item Documenta qué prompts usaste para poder reproducir tu trabajo
            \item La IA acelera el trabajo técnico, no reemplaza tu criterio
        \end{itemize}
    \end{frame}

    \begin{frame}{¿Cuándo desconfiar de un resultado?}
        \begin{itemize}
            \item Los números son "demasiado perfectos" o redondos
            \item La gráfica no cambia aunque cambies el dataset de entrada
            \item La IA no puede explicar de dónde salió un valor específico
            \item El resultado contradice algo que ya sabes del tema
        \end{itemize}
    \end{frame}

    \begin{frame}{Cierre}
        \textbf{Gracias por participar.}

        \vspace{1em}
        Todo el material de este taller está disponible en:
        \begin{itemize}
            \item \url{https://github.com/AlejandroG9/taller-ia-visualizacion-datos}
        \end{itemize}

        \vspace{1em}
        \textit{agturrubiates@uat.edu.mx}
    \end{frame}
\end{document}
